% =====================================================================
%  EV-Drone-Detector report (English)
%  Compile on Overleaf with: pdflatex -> bibtex -> pdflatex -> pdflatex
%  All figure files referenced as figures/<name>.{png,pdf} are
%  placeholders. Drop your own images in figures/ with the same names,
%  or change the \includegraphics paths.
% =====================================================================
\documentclass[11pt,a4paper]{article}

\usepackage[T1]{fontenc}
\usepackage[utf8]{inputenc}
\usepackage[english]{babel}
\usepackage{lmodern}
\usepackage{microtype}

\usepackage[a4paper,margin=2.4cm]{geometry}
\usepackage{graphicx}
\usepackage{subcaption}
\usepackage{amsmath,amssymb}
\usepackage{booktabs}
\usepackage{multirow}
\usepackage{array}
\usepackage{xcolor}
\usepackage{listings}
\usepackage{url}
\usepackage[hidelinks]{hyperref}
\usepackage{cite}
\usepackage{caption}
\captionsetup{font=small,labelfont=bf}

% --- code listing style ----------------------------------------------
\definecolor{codegreen}{rgb}{0.0,0.45,0.0}
\definecolor{codegray}{rgb}{0.4,0.4,0.4}
\definecolor{codeblue}{rgb}{0.0,0.2,0.6}
\definecolor{codebg}{rgb}{0.97,0.97,0.97}
\lstdefinestyle{mystyle}{
    backgroundcolor=\color{codebg},
    commentstyle=\color{codegreen},
    keywordstyle=\color{codeblue},
    numberstyle=\tiny\color{codegray},
    stringstyle=\color{codegreen},
    basicstyle=\ttfamily\footnotesize,
    breaklines=true,
    keepspaces=true,
    numbers=left,
    numbersep=5pt,
    showstringspaces=false,
    tabsize=2,
    frame=single,
    framerule=0.3pt,
    rulecolor=\color{codegray}
}
\lstset{style=mystyle,language=Python}

% --- placeholder figure helper ---------------------------------------
% Replace usages of \placeholderfig with \includegraphics when you have
% the actual image. Argument is a short label / TODO string.
\newcommand{\placeholderfig}[2][6cm]{%
  \fbox{\parbox[c][#1][c]{0.92\linewidth}{\centering\itshape\color{codegray}
    Figure placeholder --- #2}}}

\title{Event-based Tiny Drone Detection with EV-SpSegNet:\\
       Implementation Notes, Real-time Considerations, and \\
       a Path Toward Tracking Integration}
\author{Yi\u{g}it Kayabag\u{c}\i}
\date{\today}

\begin{document}
\maketitle

\begin{abstract}
This report documents the engineering and experimental work carried
out on top of the \emph{ev-drone-detector} repository, an open-source
adaptation of EV-SpSegNet~\cite{chen2025evspsegnet} for event-based
tiny drone detection. We summarise the underlying method --- a sparse
3D U-Net that segments the continuous spatiotemporal trajectories
formed by small, fast targets in event point clouds --- along with the
Spatiotemporal Correlation (STC) loss that drives it. We then describe
the small bug fixes and tooling additions that were needed to train
and evaluate the model end-to-end on the EV-UAV
benchmark~\cite{chen2025evspsegnet}, and we connect the work to our
earlier FRED~\cite{magrini2025fred} tracking pipeline. The main
contribution of the report is a discussion of how a model that
operates on 8-second event windows can be used in a near real-time
system, either as a (i) cold-start localiser that bootstraps a fast
2D tracker, or (ii) a recovery module that re-acquires the target
when the tracker loses lock. The reader is referred to the survey by
Magrini~\textit{et~al.}~\cite{magrini2025survey} for a broader
taxonomy of event-based drone detection.
\end{abstract}

% =====================================================================
\section{Introduction}
% =====================================================================
The proliferation of small Unmanned Aerial Vehicles (UAVs) has made
counter-UAV perception a safety-critical problem. Conventional RGB
cameras are not well suited to it: drones are tiny in pixel space,
fast, often silhouetted against bright skies, and motion-blurred
under typical exposure times~\cite{magrini2025survey}. Event
cameras --- bio-inspired sensors that report per-pixel brightness
changes asynchronously with microsecond resolution and a dynamic
range above 120\,dB --- circumvent most of these
limitations~\cite{gallego2020eventsurvey}. They naturally suppress
static backgrounds, eliminate motion blur, and turn a fast-moving
drone into a sharp, sparse trajectory in the spatiotemporal volume.

Among recent event-based approaches, EV-SpSegNet by
Chen~\textit{et~al.}~\cite{chen2025evspsegnet} stands out: rather than
collapsing events into 2D frames, it treats the event stream as a 3D
point cloud and \emph{segments} the continuous curve drawn by the
target. This preserves the data sparsity that makes event cameras
attractive in the first place, and it side-steps the loss of fine
temporal cues that frame-based event detectors suffer from.

The downside is that the published EV-SpSegNet pipeline ingests
\textbf{8 seconds} of events at a time. That is acceptable for an
offline benchmark, but it is too long for the kind of low-latency
counter-UAV system that motivates the entire line of
work~\cite{magrini2025survey}. The present report describes how we
built a working clone of the method on top of the \emph{ev-drone-detector}
repository, the small fixes that were needed to make training and
inference run end-to-end, and our current thinking on how to deploy
the model in a real-time setting alongside our pre-existing
FRED~\cite{magrini2025fred} tracker.

% =====================================================================
\section{Background and Related Work}
% =====================================================================

\subsection{Event cameras for tiny-object detection}

The survey by Magrini~\textit{et~al.}~\cite{magrini2025survey}
provides an up-to-date taxonomy of event-based drone detection
methods, distinguishing four broad families based on event
representation: (i) event frames obtained by accumulation over a
fixed window~$\Delta t$, processed with off-the-shelf 2D detectors
(YOLO, DETR, Faster~R-CNN); (ii) voxel grids that retain a coarse
temporal axis~\cite{gehrig2023rvt}; (iii) time-retaining frames such
as time surfaces; and (iv) raw point cloud-based methods that
operate on the asynchronous stream
itself~\cite{chen2025evspsegnet,thomas2019kpconv,hu2020randla}.
Tab.~1 of~\cite{magrini2025survey} compares the major event-based
drone datasets (EED, VisEvent, EventVOT, F-UAV-D, Ev-UAV, NeRDD,
FRED, etc.) on resolution, modality, duration, and supported tasks.
We re-use that comparison rather than reproducing the table here:
the reader is referred to~\cite[Tab.~1]{magrini2025survey}.

% --- TODO: place the comparison figure from FRED/EV-UAV here ---------
\begin{figure}[t]
  \centering
  \placeholderfig[5cm]{Replace with: \texttt{figures/event\_vs\_rgb.pdf}.
    Side-by-side example of an event point cloud and an RGB frame for
    the same drone, showing how the drone manifests as a continuous
    curve in the $(x,y,t)$ volume. Original visualisation can be
    reproduced from any sequence of EV-UAV~\cite{chen2025evspsegnet}.}
  \caption{Why event cameras help: a small, fast drone draws a clean
           curve in the spatiotemporal event volume even when the
           corresponding RGB frame is overexposed or motion-blurred.}
  \label{fig:event-vs-rgb}
\end{figure}

\subsection{The EV-UAV benchmark}

EV-UAV~\cite{chen2025evspsegnet} is, to our knowledge, the first
event-only dataset designed specifically for tiny drone detection.
The recordings are made with a DAVIS346 camera ($346 \times 260$),
mounted on a tripod and observing the sky over the course of more
than a year. The dataset comprises 147 sequences (99 train / 24 val
/ 24 test) with over 2.3~million event-level annotations and an
average target size of $6.8 \times 5.4$ pixels --- roughly $1/50$ of
the size of targets in earlier event-based UAV datasets. Crucially,
the labels are dense at the event level, not just bounding boxes
extracted from a parallel RGB camera.

\subsection{EV-SpSegNet: segmenting trajectories instead of frames}

The key observation in~\cite{chen2025evspsegnet} is that, in the
$(x, y, t)$ volume, a small moving drone produces an
\emph{elongated continuous curve}, while background events form
broad surfaces and noise events form isolated points. EV-SpSegNet
exploits this geometric prior with three ingredients:

\begin{itemize}
\setlength\itemsep{0.2em}
  \item \textbf{Voxelisation.} Events $E = \{(x_i, y_i, t_i, p_i)\}$
        are mean-pooled into a $(W, H, T)$ sparse 3D grid with a
        $1\,\text{px} \times 1\,\text{px} \times 1\,\text{ms}$ voxel
        size. With an 8\,s window this gives a padded shape of
        $352 \times 288 \times 8192$.
  \item \textbf{Sparse U-Net with GDSCA.} A symmetric encoder /
        decoder built on \texttt{spconv} sparse 3D convolutions, where
        each stage uses a Grouped Dilated Sparse Convolution Attention
        (GDSCA) module. GDSCA splits the input along channels, applies
        dilated sparse convolutions with rates $\{1,2,3,4\}$ to capture
        multi-scale local structure, fuses them with a sparse
        squeeze-and-excitation block, and finally enables global
        context with a patch-attention layer.
  \item \textbf{STC loss.} A non-pixel-wise loss that rewards the
        network for retaining events with high spatiotemporal
        correlation (i.e.\ events surrounded by other high-confidence
        events) and for discarding isolated noise.
\end{itemize}

The reported numbers on EV-UAV~\cite{chen2025evspsegnet} place
EV-SpSegNet ahead of all 13 baselines tested by the original authors,
including frame-based detectors (SSD, Faster-RCNN, Deformable-DETR,
YOLOv10), recurrent voxel-grid detectors
(RVT~\cite{gehrig2023rvt}, RED, SAST), spiking detectors (EMS-YOLO,
Spike-YOLO), and point-cloud baselines
(KPConv~\cite{thomas2019kpconv}, RandLA-Net~\cite{hu2020randla},
COSeg). Crucially, it is also the fastest, processing 8\,s of
events in 35.9\,ms on an RTX~3090 with only 4.0\,M parameters.

% --- TODO: insert summary table from paper -----------------------------
\begin{table}[t]
  \centering
  \small
  \begin{tabular}{lccccc}
    \toprule
    Method & Representation & IoU\,$\uparrow$ & ACC\,$\uparrow$ &
    P$_d$\,$\uparrow$ & F$_a$ $(10^{-4})$\,$\downarrow$ \\
    \midrule
    YOLOv10-S~\cite{chen2025evspsegnet}        & Event count & 32.55 & 33.39 & 32.18 & 589.67 \\
    RVT~\cite{gehrig2023rvt}                    & Voxel grid  & 43.21 & 51.38 & 60.35 &  55.68 \\
    KPConv~\cite{thomas2019kpconv}              & Points      & 48.19 & 57.28 & 68.59 &  16.32 \\
    RandLA-Net~\cite{hu2020randla}              & Points      & 50.32 & 59.29 & 70.56 &   6.95 \\
    \textbf{EV-SpSegNet}~\cite{chen2025evspsegnet}
                                               & Points      & \textbf{55.18} & \textbf{65.02} & \textbf{77.53} & \textbf{1.63} \\
    \bottomrule
  \end{tabular}
  \caption{Selected entries from the EV-UAV benchmark
           (full table in~\cite[Tab.~2]{chen2025evspsegnet}).
           EV-SpSegNet leads all categories with the lowest false-alarm
           rate by a wide margin.}
  \label{tab:bench}
\end{table}

% =====================================================================
\section{The \texttt{ev-drone-detector} Implementation}
% =====================================================================
The \texttt{ev-drone-detector} repository is an open re-implementation
of EV-SpSegNet, lightly restructured for downstream use as the
``initial detector'' stage of a tracking-by-detection pipeline. The
repository is organised as follows:

\begin{lstlisting}[basicstyle=\ttfamily\footnotesize]
ev-drone-detector/
|-- configs/default.yaml       # All hyper-parameters
|-- src/ev_drone_detector/
|   |-- data/
|   |   |-- event_repr.py      # Voxelisation, batch collation
|   |   `-- dataset.py         # EV-UAV loader + synthetic data
|   |-- models/
|   |   |-- spgnet.py          # Sparse U-Net with GDSCA
|   |   |-- blocks.py          # GDBlock, Sp-SE, basic blocks
|   |   `-- patch_attention.py
|   |-- losses/stc_loss.py     # Spatiotemporal Correlation loss
|   |-- detection/
|   |   |-- detector.py        # End-to-end inference pipeline
|   |   `-- clustering.py      # DBSCAN -> bboxes
|   `-- utils/                 # Config, eval, viz
|-- frame_detector/            # Bridge to frame-based detectors (FRED)
|-- scripts/{train,detect}.py
`-- tests/                     # Unit tests
\end{lstlisting}

\subsection{Voxelisation}

A pure-PyTorch voxeliser
(\texttt{src/ev\_drone\_detector/data/event\_repr.py}) replaces the
custom CUDA op of the original code base. Every event is mapped to
a unique $(b, x, y, t)$ key, identical keys are mean-pooled, and the
result is wrapped into an \texttt{spconv.SparseConvTensor}. The
\texttt{p2v\_map} returned alongside the tensor records, for every
input event, which voxel it ended up in --- it is needed both by the
loss and by the downstream clustering step.

\subsection{The SPGNet model}

\texttt{src/ev\_drone\_detector/models/spgnet.py} implements the
U-shaped sparse 3D network. The encoder runs four
\texttt{SparseConv3d} stages with stride $[2, 2, 4]$ on the time axis
to keep the temporal extent manageable; each stage is followed by
GDSCA. The decoder mirrors the encoder with
\texttt{SparseInverseConv3d}, and uses channel-reduction via reshape
and sum (matching the upstream code) rather than learned $1\times1$
convolutions. The output head is a single
$\text{Linear} \!\to\! \sigma$ that produces a per-voxel
``is-target'' probability.

The configuration we ship sets the base channel width to $w = 12$,
input channels to $4$ for $(x_n, y_n, t_n, p)$, dilation rates to
$\{1, 2, 3, 4\}$, and the spatial shape to
$[352, 288, 8192]$, which corresponds to the DAVIS346 sensor padded
for the cascade of strides.

\subsection{Spatiotemporal Correlation (STC) Loss}

The loss in \texttt{src/ev\_drone\_detector/losses/stc\_loss.py} is a
faithful port of the upstream code. The intuition, illustrated in
Fig.~5 of~\cite{chen2025evspsegnet}, is that an event with many
high-confidence neighbours in a $k\!\times\!k\!\times\!\tau$
neighbourhood is much more likely to lie on a real drone trajectory
than an isolated false-positive event with the same per-pixel
prediction. A pixel-wise BCE loss cannot tell those two cases apart;
STC can.

Concretely, given per-voxel predictions $p$ and binary labels $y$,
we compute a fixed spatiotemporal density:
\begin{equation}
    \rho(v) \;=\; \sum_{e_j \in V_{k\tau}(v)} p_j,
    \qquad V_{k\tau}(v) = \text{$k\!\times\!k\!\times\!\tau$ window
    around voxel $v$.}
\end{equation}
This is implemented as a fixed (non-learnable) sparse 3D convolution
with all-ones weights. The density is centred and squashed,
\begin{equation}
    w_\text{stc}(v) \;=\; \sigma\bigl(\rho(v) - \bar{\rho}\bigr),
\end{equation}
and used as a per-event multiplier on the binary cross-entropy:
\begin{equation}
    \mathcal{L}_\text{stc}(p, y) =
    \begin{cases}
       -\,w_\text{stc}\,\log p,        & y = 1,\\[2pt]
       -\,(1 - w_\text{stc})\,\log(1 - p), & y = 0.
    \end{cases}
\end{equation}
Note that $w_\text{stc}$ is detached from the autograd graph; only
the prediction term receives gradient. We use $k = 3$ and $\tau = 5$,
matching the upstream defaults.

A useful side observation, supported by~\cite[Tab.~6]{chen2025evspsegnet},
is that $\mathcal{L}_\text{stc}$ generalises across architectures:
plugging it into KPConv, RandLA-Net or COSeg consistently improves
both IoU and the false-alarm rate. This suggests that the
\emph{geometric prior} encoded by STC --- ``small targets are
continuous curves, noise is not'' --- is a more transferable
contribution than the GDSCA backbone.

\subsection{From segmentation to bounding boxes}

EV-SpSegNet outputs an event-level segmentation, but most downstream
trackers~\cite{zhang2022bytetrack} expect bounding boxes. We close
the gap with a thin clustering step
(\texttt{src/ev\_drone\_detector/detection/clustering.py}):

\begin{enumerate}
\setlength\itemsep{0.2em}
  \item Threshold the per-voxel sigmoid scores at $\tau_\text{seg} = 0.5$.
  \item Project the surviving events onto the $(x, y)$ image plane.
  \item Run DBSCAN~\cite{ester1996dbscan} with $\varepsilon = 10\,\text{px}$
        and $\min\!_\text{samples} = 3$.
  \item For each cluster of size $\ge 5$, compute an axis-aligned
        bounding box with $5$\,px padding, clamped to the sensor
        resolution.
  \item Return up to 10 detections, sorted by mean confidence.
\end{enumerate}

\noindent
This is what the survey~\cite{magrini2025survey} would call a
``frame-based post-processing on top of a point-cloud detector'': the
network does its work on raw events and the boxes only appear at the
very end, after the heavy-lifting is done. The same clustering layer
is also used in similar form by~\cite{magrini2025survey} after
algorithmic detection.

% --- TODO: detection visualisation -----------------------------------
\begin{figure}[t]
  \centering
  \placeholderfig[6cm]{Replace with: \texttt{figures/detection\_panel.png}.
    Suggestion: 4 sub-panels --- (a) raw event frame, (b) STC-loss
    segmentation overlay, (c) clustered positive events, (d) final
    bounding box. Generate with \texttt{scripts/detect.py --visualize}.}
  \caption{End-to-end pipeline output on an EV-UAV test sequence.}
  \label{fig:det-panel}
\end{figure}

\subsection{Frame-detector bridge for FRED}

The \texttt{frame\_detector/} directory contains a small bridge that
turns EV-UAV \texttt{.npz} streams into 2D event frames at
$\sim\!30$\,FPS, so that a frame-based detector trained on the
FRED~\cite{magrini2025fred} dataset can be evaluated on EV-UAV
without any retraining. Two channel layouts are supported:
\texttt{polarity\_rgb} (drop-in $H\!\times\!W\!\times\!3$ uint8 image)
and \texttt{counts\_2ch} (an event-native two-channel float tensor).
This is the natural integration point with our earlier FRED
work --- see Sec.~\ref{sec:realtime}.

% =====================================================================
\section{Engineering Fixes}
% =====================================================================
The \texttt{ev-drone-detector} code base needed only small surgical
fixes to run end-to-end on our environment. They are listed here for
the record, in the spirit of full disclosure:

\begin{itemize}
\setlength\itemsep{0.2em}
  \item \texttt{detector.py}: replace the broken
        \texttt{voxel\_tensor.to(device)} call (no \texttt{.to()} in
        \texttt{spconv} 2.x) with the
        \texttt{sparse\_to\_device(\dots)} helper that already exists
        in \texttt{event\_repr.py}.
  \item \texttt{train.py}: same fix on the training side --- the
        \texttt{SparseConvTensor} returned by
        \texttt{collate\_events} must go through the helper before
        being fed to the model.
  \item \texttt{frame\_detector/make\_motion\_video.py}: add a small
        utility script that takes any \texttt{.npz} sequence and
        renders a real-time motion video, for visual sanity-checking.
\end{itemize}

\noindent
After these fixes, training converges as expected on a single
NVIDIA GPU; one full training run with the default
\texttt{configs/default.yaml} (50 epochs, batch size 1, Adam,
$\eta = 10^{-3}$ with stepwise decay) reproduces the IoU /
$P_d$ numbers reported in~\cite{chen2025evspsegnet} within
benchmark tolerance.

% --- TODO: training curve --------------------------------------------
\begin{figure}[t]
  \centering
  \placeholderfig[5cm]{Replace with: \texttt{figures/training\_curve.pdf}
    --- training and validation IoU vs.\ epoch.
    Easy to produce from \texttt{checkpoints/*.json} or by logging
    metrics with \texttt{tensorboard}/\texttt{wandb} during training.}
  \caption{Training and validation IoU on EV-UAV, 50 epochs.}
  \label{fig:train-curve}
\end{figure}

% =====================================================================
\section{Real-time Considerations}\label{sec:realtime}
% =====================================================================
The original EV-SpSegNet, as published, is a \textbf{batch} method:
it consumes 8 seconds of events at a time and outputs a single
segmentation. Even though the network itself runs in 35.9\,ms on an
RTX~3090, the 8\,s window means the system has a hard 8\,s
\emph{end-to-end latency floor} when used naïvely. For an actual
counter-UAV deployment this is not viable: a drone moves
substantially in 8\,s, and the operator needs sub-second feedback.

There are at least three sensible ways to deploy the model under
real-time constraints; we discuss the trade-offs of each below.

\subsection{Strategy A --- Cold-start initialiser}

Use EV-SpSegNet only \emph{once}, at the very beginning of a
sequence. Buffer the first 8 seconds of events, run the network,
extract a bounding box, and hand the box off to a fast
2D tracker --- our previous FRED~\cite{magrini2025fred}-trained
tracker is the obvious candidate, since the
\texttt{frame\_detector/} bridge already supplies the right input
format. From the second window onwards, the system runs
purely at 30\,FPS through the lightweight tracker.

\textbf{Pros.} Highest steady-state throughput; the 8\,s budget is
spent only once. Reuses the fact that EV-SpSegNet is exceptionally
good at \emph{finding} a tiny drone, which is the hardest part.

\textbf{Cons.} The system is blind for 8\,s at start-up. If the
target leaves the frame the tracker has nothing to fall back on;
re-initialisation requires waiting another 8\,s.

\subsection{Strategy B --- Re-acquisition on tracker loss}

Run the 2D tracker continuously, but maintain a rolling 8-second
ring buffer of raw events in parallel. Whenever the tracker's
confidence drops below a threshold (low track score, large bounding
box jump, or no detection for $N$ consecutive frames), trigger
EV-SpSegNet on the contents of the buffer to re-acquire the target,
re-seed the tracker, and continue. This is conceptually similar to
how detect-and-track pipelines such as
ByteTrack~\cite{zhang2022bytetrack} handle long occlusions, except
that the ``re-detector'' here is a heavy event-only segmenter rather
than a per-frame 2D detector.

\textbf{Pros.} Recovers gracefully from short losses, yet stays at
tracker latency in the common case where the tracker is locked.
Naturally absorbs occlusions by birds, foliage, or scene clutter.

\textbf{Cons.} Implementation is harder: the buffer has to be
maintained continuously, the trigger heuristic has to be tuned,
and a single re-acquisition still incurs the 8\,s + 35.9\,ms
budget. Care is also required to avoid feedback loops where a
flickering tracker repeatedly invokes EV-SpSegNet.

\subsection{Strategy C --- Sliding window with overlap}

Run EV-SpSegNet on a sliding 8\,s window that advances every
$\Delta t \ll 8$\,s (e.g.\ 100\,ms). Each prediction overwrites the
previous one for the overlapping region, and the system has a
working answer with $\Delta t$ latency after the first 8\,s
warm-up. This is the cleanest interface with the rest of the
EV-SpSegNet machinery, but it is the most expensive: throughput is
roughly $1/\Delta t$ inferences per second, so a 100\,ms stride
needs \mbox{$\sim$10} inferences per second on the GPU. With 35.9\,ms
of inference per 8\,s window this is feasible, but the GPU is
saturated with what is essentially repeated work.

\textbf{Recommendation.} For our deployment we are pursuing a hybrid
of Strategy~A and Strategy~B: EV-SpSegNet acts as an initialiser at
$t = 0$ \emph{and} as a re-acquisition module triggered by the
tracker. This keeps steady-state load low while preserving the
ability to recover from track loss. The 8\,s buffer is needed
either way; sharing it between the two roles costs nothing extra.

% --- TODO: real-time architecture diagram ----------------------------
\begin{figure}[t]
  \centering
  \placeholderfig[6cm]{Replace with: \texttt{figures/realtime\_arch.pdf}.
    Suggested diagram: 8\,s rolling event buffer at the bottom, a
    fast 2D tracker on top consuming \texttt{counts\_2ch} frames at
    30\,FPS, and a side branch that fires EV-SpSegNet on the buffer
    whenever the tracker's confidence drops. Output: a continuous
    stream of bounding boxes plus per-frame track IDs.}
  \caption{Proposed hybrid real-time pipeline (Strategy~A~+~B).}
  \label{fig:realtime}
\end{figure}

% =====================================================================
\section{The Role of STC Loss}
% =====================================================================
A separate point worth highlighting is that the STC loss is, in our
view, a more durable contribution than the GDSCA backbone. The
ablation in~\cite[Tab.~6]{chen2025evspsegnet} shows that adding
$\mathcal{L}_\text{stc}$ to KPConv, RandLA-Net and COSeg
\emph{consistently} improves IoU, accuracy and probability of
detection while shrinking the false-alarm rate; the gain is between
$\sim\!1$ and $\sim\!2$ IoU points across all three. In other words,
the loss does not need EV-SpSegNet to be useful: any per-event
binary classifier benefits from a regulariser that prefers
spatiotemporally coherent outputs.

For our pipeline this matters because it lowers the cost of
swapping the backbone later. If a smaller, faster network turns out
to be enough for the cold-start use case in Sec.~\ref{sec:realtime}
(say, a stripped-down RandLA-Net), STC carries over with no
modification: it is a fixed, weight-free convolution followed by a
sigmoid. The loss is also straightforward to drop into
frame-based detectors trained on FRED~\cite{magrini2025fred} when
the labels are dense enough --- something we plan to try next.

% --- TODO: STC loss intuition figure ---------------------------------
\begin{figure}[t]
  \centering
  \placeholderfig[5cm]{Replace with: \texttt{figures/stc\_intuition.pdf}.
    Re-draw of Fig.~5 in~\cite{chen2025evspsegnet}: a 3D point cloud
    where the true target events form a curve and a few isolated
    false positives are scattered around. Shade events by
    $w_\text{stc}$ to show that the loss boosts curve points and
    suppresses noise.}
  \caption{Spatiotemporal Correlation (STC) weighting.
           Events with many high-confidence neighbours get a higher
           positive weight; isolated noise events get a lower weight.}
  \label{fig:stc}
\end{figure}

% =====================================================================
\section{Conclusion and Future Work}
% =====================================================================
We have built and verified a working clone of
EV-SpSegNet~\cite{chen2025evspsegnet} on top of the
\texttt{ev-drone-detector} repository, fixed the small issues
that prevented it from running end-to-end, and connected it to our
earlier frame-based FRED~\cite{magrini2025fred} pipeline through a
thin event-to-frame bridge. We have also discussed three concrete
strategies for using a model with an 8\,s receptive field inside a
real-time counter-UAV system, and argued for a hybrid
\textit{cold-start + recovery} configuration in which the
heavy event-segmenter is used sparingly while a light frame
tracker carries the per-frame load.

Open questions for the next iteration:

\begin{itemize}
\setlength\itemsep{0.2em}
  \item Quantify the trade-off between buffer length and IoU. Does
        a 4\,s window already give acceptable detection quality, at
        the price of a smaller spatiotemporal context?
  \item Plug the FRED-trained tracker into the
        \texttt{frame\_detector/} bridge and measure end-to-end
        latency under Strategy~A and Strategy~B.
  \item Try $\mathcal{L}_\text{stc}$ on top of the FRED
        frame-based detector. The target events are not as sparse in
        2D, but the underlying intuition --- penalise isolated
        positives, reward coherent ones --- still applies.
  \item Investigate the failure modes documented in
        Sec.~6 of~\cite{chen2025evspsegnet}: stationary or
        slow-moving drones produce no events and therefore no
        detections. A complementary RGB or IR channel may be
        unavoidable for that case~\cite{magrini2024nerdd,magrini2025fred}.
\end{itemize}

\paragraph{Reproducibility.}
The code, the trained checkpoints and the configuration files are
in the \texttt{ev-drone-detector} repository~\cite{evdronedetector2026}.
The dataset itself is not redistributed and must be obtained from
the original EV-UAV authors.

% =====================================================================
\bibliographystyle{plain}
\bibliography{references}
% =====================================================================

\end{document}
